
\section{Benchmark implementation details}
\label{app:bench}

All four benchmarks share the same loop (Sec.~\ref{sec:loop}) and the same
acceptance test: a proposal is accepted if its canonical string equals the
canonical string of the target's action on the same context. Canonical strings
are produced by the benchmark's parser from the raw response; a response that
does not parse (no action, several actions where one is required, an unknown
tool, or arguments that violate the tool schema) never matches. The target
decodes greedily (temperature 0, seed 0) in every end-to-end run. Table
\ref{tab:bench} lists the per-benchmark settings; the paragraphs below give the
formats.

\begin{table}[h]
\centering
\footnotesize
\setlength{\tabcolsep}{4pt}
\caption{Per-benchmark settings of the end-to-end runs.}
\label{tab:bench}
\resizebox{\linewidth}{!}{%
\begin{tabular}{>{\raggedright\arraybackslash}p{2.5cm}>{\raggedright\arraybackslash}p{2.9cm}>{\raggedright\arraybackslash}p{3.1cm}>{\raggedright\arraybackslash}p{3.0cm}>{\raggedright\arraybackslash}p{3.3cm}}
\toprule
 & OpenAGI & TaskBench & $\tau^2$-bench & BFCL multi-turn \\
\midrule
task type & tool-selection plan & tool calls with args & ticket-driven support & multi-turn function calling \\
tools & 14 $+$ TERMINATE & DailyLife APIs & 44 & per-task catalog \\
target format & ReAct, \texttt{<tool>} tag with tool name & ReAct, \texttt{<tool>} tag with name(\{json\}) & single-call JSON, guided & \texttt{<tool\_call>} tags, ordered group \\
target reasoning & thought $\le 512$ tok. & thought $\le 512$ tok. & on, budget 2{,}048 & on, budget 4{,}096 \\
target action budget & 32 tok. & 128 tok. & (in budget) & (in budget) \\
draft format & direct, tool name & direct, name$+$args & grammar, JSON call & grammar, call group / \texttt{<respond/>} \\
draft budget & 32 tok. & 128 tok. & 512 tok. & 512 tok. \\
observation & none (plan-first) & none (plan-first) & tool result (local backend) & tool result (cloned runtime) \\
segment boundary & task end & task end & task end (\texttt{done}) & user turn \\
$k$ & 3 & 3 & 2 & 2 \\
max steps & 12 & 12 & 60 & 20 per turn \\
train / test tasks & 250 / 63 & 228 / 60 & 60 (+14 val) / 40 & 130 / 60 \\
training pairs & 1{,}090 & 1{,}120 & 932 & 1{,}254 \\
\bottomrule
\end{tabular}}
\end{table}

\paragraph{OpenAGI.}
The agent selects, one step at a time, which of 14 domain-expert tools to
apply to a task description~\citep{openagi}, and emits \texttt{TERMINATE} when
the plan is complete. The target runs ReAct: a thought of at most 512 tokens
followed by an action of at most 32 tokens in the form
\texttt{<tool>Image Denoising</tool>}; the parser takes the last tag pair so
that a preamble naming other tools does not win over the committed action. The
canonical action is the tool name. Steps do not observe tool outputs, so the
context of step $i$ is the task and the committed actions $a_{<i}$ (the target's
thoughts are not part of the context and are not shown to the draft). A task
ends at \texttt{TERMINATE} or after 12 steps. The 313 tasks are split by task
into 250 training and 63 test tasks (seed 42); running the target once over the
training tasks yields 1{,}090 (context, action) pairs, and the test tasks give
263 held-out contexts for the offline agreement in Table~\ref{tab:offline}. The
draft receives the same context rendered as a direct prompt and answers with
the action tag only.

\paragraph{TaskBench.}
We use the DailyLife API subset of TaskBench~\citep{taskbench}, where each
step is a tool call with JSON
arguments, or \texttt{TERMINATE}; for example\par\noindent\hspace*{1.5em}%
\texttt{<tool>apply\_for\_passport(\{"country":"Germany"\})</tool>}\par\noindent The 300 tasks were split 240 / 60 by task before any
model call. The target runs ReAct as on OpenAGI (thought $\le 512$ tokens, action
$\le 128$ tokens). The parser is strict: the response must contain exactly one
action, the tool must exist, the argument names must equal the schema's
parameter set, and every argument value must be a string; anything else is a
parse or schema failure and cannot match. The canonical action is
\texttt{name(\{sorted-key JSON\})}. The training set is derived
deterministically from the target's raw responses on the training tasks
(228 tasks with a valid trajectory, 1{,}120 pairs; 1{,}011 / 109 split by task
for validation); the same derivation on the test tasks gives 57 tasks / 277
contexts for offline agreement, and all 60 test tasks are run end to end. No
test task's tool sequence occurs in training.

\paragraph{$\tau^2$-bench (telecom domain).}
We use the telecom domain of $\tau^2$-bench~\citep{tau2bench} in its
\emph{solo} configuration: there is no simulated user; the agent works from the
task ticket, a policy document, and 44 tools whose implementations run against
the benchmark's local state. We compact the policy text (same rules, shorter
wording) so that contexts fit the serving window; the compact policy is
identical for both arms and both models. The target receives the tool schemas
natively and must emit exactly one call per step as a JSON object
\texttt{\{"name":\,\ldots,"arguments":\,\ldots\}}, constrained by a
\texttt{oneOf} schema over the 44 tools (guided decoding), with reasoning
enabled and a 2{,}048-token budget; a task ends with the \texttt{done} call or
after 60 steps, and success is the benchmark's own reward function
(\texttt{evaluate\_simulation}, solo mode). The canonical action is the JSON
object with sorted keys. Training and validation pairs were collected from the
target at its default sampling for this mode (temperature 0.6, top-$p$ 0.95,
top-$k$ 20, seed 0); end-to-end runs use temperature 0. The draft sees the same
context rendered as a direct prompt and decodes a \texttt{<tool>\{json\}</tool>}
block under a grammar derived from the same schemas (512-token budget, no
reasoning). Each proposal is executed in a cloned environment branch so that
the next proposal sees its observation; a rejected branch is discarded.

\paragraph{BFCL multi-turn.}
We use the multi-turn split of the Berkeley Function Calling
Leaderboard~\citep{bfcl}, 130 training and 60 test tasks.
Each task is a conversation; in each user turn the agent issues one or more
ordered groups of function calls and ends the turn with a natural-language
reply. The target emits calls as \texttt{<tool\_call>\{json\}</tool\_call>}
tags (several tags form one ordered group, which is one decision), with
reasoning enabled, a 4{,}096-token budget, and at most 20 decisions per turn.
The draft predicts the next decision as
\texttt{<tool>[\{call\},\ldots]</tool>} or the literal \texttt{<respond/>},
decoding under a GBNF grammar built from the tools visible in the context
(512-token budget, no reasoning); it never predicts reply text. A decision
matches when the ordered call group is identical after canonicalization (JSON
with sorted keys) or when both sides end the turn. Proposals execute in a
cloned tool runtime (deep-copied instances); the committed branch's state is
the one the next turn sees. User turns are segment boundaries: at the end of
each turn a width-1 probe lets the target confirm the turn end. Official
success is computed by the official BFCL evaluator on the final state and
responses. Verification uses prefix-ordered cancellation
(App.~\ref{app:serving}); the 60 test tasks were run twice end to end, one task
at a time. On the 10-task development split, run twice at each depth before
the test runs, the speculative arm was faster at $k{=}1$ and $k{=}2$
(1.14/1.23\xx{} and 1.18/1.24\xx{}) but completed fewer tasks (9 against 14 at
$k{=}1$, 13 against 14 at $k{=}2$); we selected $k{=}2$ and report the test
split, where success was higher in both passes (Table~\ref{tab:four}).

\section{Trace-replay predictor and idle-time accounting}
\label{app:predictor}

\paragraph{Inputs.}
The predictor of Sec.~\ref{sec:validation} uses (i) the sequential arm's
per-step generation times, grouped by segment; (ii) the per-position hit rates
$q_j$; and (iii) the mean draft time per round from the speculative arm's draft
records. It uses no timing from the speculative arm's verifications.

\paragraph{Replay.}
For every segment of every task, the replay walks the sequential steps in
order. At the start of a round with $r$ steps remaining it draws the number of
leading hits $A$ from independent Bernoulli trials with rate $q_1$, capped at
$\min(k, r-1)$; the round commits $A{+}1$ steps and its wait is the maximum of
those steps' sequential generation times. The predicted $\beta$ is the sum of
round waits divided by the number of rounds times $t_{\mathrm{gen}}$, averaged
over 50 replays. The commits per round $c$ in Table~\ref{tab:model} is the
measured value (Eq.~\ref{eq:c} with the measured $q_j$ and $e$); the
i.i.d.\ column uses the replay's own $A$ instead. Table~\ref{tab:qj} lists the
measured $q_j$ and $e$ for the nine configurations.

\begin{table}[h]
\centering
\small
\caption{Measured per-position hit rates $q_j = \Pr[A \ge j]$, end-of-segment
loss $e$, and committed steps per round $c = 1 + \sum_j q_j - e$. A dash marks a position beyond the row's $k$.}
\label{tab:qj}
\begin{tabular}{llrrrrrrrr}
\toprule
benchmark & draft & $k$ & $q_1$ & $q_2$ & $q_3$ & $q_4$ & $q_5$ & $e$ & $c$ \\
\midrule
OpenAGI   & trained 0.6B   & 3 & 0.71 & 0.44 & 0.28 & --   & --   & 0.38 & 2.05 \\
OpenAGI   & trained 0.6B   & 1 & 0.74 & --   & --   & --   & --   & 0.20 & 1.54 \\
OpenAGI   & trained 0.6B   & 5 & 0.70 & 0.47 & 0.31 & 0.14 & 0.08 & 0.46 & 2.25 \\
OpenAGI   & untrained 0.6B & 3 & 0.15 & 0.00 & 0.00 & --   & --   & 0.01 & 1.14 \\
OpenAGI   & 32B-direct     & 3 & 0.70 & 0.50 & 0.33 & --   & --   & 0.34 & 2.19 \\
TaskBench & trained 0.6B   & 3 & 0.66 & 0.34 & 0.18 & --   & --   & 0.30 & 1.88 \\
TaskBench & 32B-direct     & 3 & 0.74 & 0.50 & 0.33 & --   & --   & 0.36 & 2.21 \\
$\tau^2$-bench   & trained 0.6B   & 2 & 0.64 & 0.35 & --   & --   & --   & 0.11 & 1.88 \\
BFCL      & trained 0.6B   & 2 & 0.38 & 0.07 & --   & --   & --   & 0.00 & 1.45 \\
\midrule
\multicolumn{10}{l}{\emph{Draft-size study, common serving configuration (Sec.~\ref{sec:training})}} \\
OpenAGI   & trained 0.6B   & 3 & 0.69 & 0.44 & 0.27 & --   & --   & 0.37 & 2.04 \\
OpenAGI   & trained 1.7B   & 3 & 0.74 & 0.50 & 0.33 & --   & --   & 0.37 & 2.20 \\
OpenAGI   & trained 4B     & 3 & 0.79 & 0.50 & 0.31 & --   & --   & 0.34 & 2.26 \\
TaskBench & trained 0.6B   & 3 & 0.67 & 0.37 & 0.20 & --   & --   & 0.30 & 1.94 \\
TaskBench & trained 1.7B   & 3 & 0.69 & 0.42 & 0.28 & --   & --   & 0.34 & 2.05 \\
TaskBench & trained 4B     & 3 & 0.71 & 0.43 & 0.27 & --   & --   & 0.31 & 2.10 \\
\bottomrule
\end{tabular}
\end{table}

\paragraph{Tool terms.}
With a delay $\delta$ added to every executed tool step and $\tau = \delta / t_{\mathrm{gen}}$, the sequential
arm pays $\tau$ (in target steps) on every step that executes a tool. A round pays it for the proposals' tools,
which run while proposing and are part of the measured draft term $D$, and for every tool step the target
itself executes. Eq.~\ref{eq:S} becomes $S = c\,(1 + f\tau) / (\beta + D + g\tau)$, with $f$ the share of
sequential steps that execute a tool and $g$ the number of tool steps per round that the target executes. On
$\tau^2$-bench every step but the final one executes a tool and each round ends with one target action, and we
use $f = g = 1$. On BFCL, turn ends execute no tool ($f = 0.62$), accepted proposals are adopted rather than
executed again, so only the target's corrections and bonus actions count ($g = 0.44$--$0.45$), and $D$ is the
rollout time per round, which contains the proposals' tool steps.

\paragraph{Decomposition of the speculative arm.}
For each round we record the wall-clock of every target verification and the
round's total verification wall-clock. The speculative arm's task time then
decomposes as
\[
  \text{spec} \;=\; \sum_{\text{rounds}} \max_{\text{used}} \text{wall}
  \;+\; \text{idle} \;+\; \text{draft} \;+\; \text{residual},
\]
where ``used'' are the verifications whose result the round consumed, idle is
$\sum_{\text{rounds}} (\text{verification wall} - \max_{\text{used}}
\text{wall})$, i.e.\ time spent waiting for verifications that were already
discarded, draft is the summed draft time, and residual covers tool execution
and bookkeeping. With prefix-ordered cancellation the idle time is 0\,s of
199\,s on OpenAGI, 0\,s of 400\,s on TaskBench, 3\,s of 4{,}367\,s on $\tau^2$-bench
and 14\,s of 11{,}132\,s on BFCL, and at most 4\,s in every other row of Table~\ref{tab:model}.
The earlier gather-mode runs waited for discarded verifications: 9\,s of 205\,s (OpenAGI), 12\,s of
391\,s (TaskBench) and 552\,s of 4{,}882\,s ($\tau^2$-bench) for the trained draft, 9\,s of 267\,s and 16\,s
of 481\,s for 32B-direct on OpenAGI and TaskBench, and 242\,s of 544\,s for the untrained draft.
Re-running every row with cancellation moved the main speedups from 1.46, 1.37 and 1.41\xx{} to 1.51, 1.37
and 1.58\xx{}, the untrained draft from 0.55\xx{} to 0.98\xx{}, 32B-direct from 1.12 and 1.11\xx{} to 1.16
and 1.13\xx{}, and OpenAGI at $k{=}1$ and $k{=}5$ from 1.30 and 1.48\xx{} to 1.34 and 1.57\xx{}. Residual
is at most 83\,s (BFCL). ``No idle'' in Table~\ref{tab:model} is
$\text{seq} / (\text{spec} - \text{idle})$.

\paragraph{Concurrency cost.}
To check assumption (b) of Sec.~\ref{sec:model} we compare milliseconds per
generated token in the sequential arm with the same quantity over the used
verifications of the speculative arm: 25.5 vs.\ 26.7 (OpenAGI), 25.5 vs.\ 26.3
(TaskBench), 26.0 vs.\ 26.7 ($\tau^2$-bench), 26.1 vs.\ 26.4 (BFCL). Tokens per step
are unchanged on three benchmarks and 5\% lower on $\tau^2$-bench.

\section{Serving configuration}
\label{app:serving}

All runs use vLLM on one NVIDIA H100 (80\,GB). The target is
Qwen3-32B (bf16, tensor parallel 1) and the draft is Qwen3-0.6B (bf16) with a
LoRA adapter of rank 64
served through vLLM's LoRA support (\texttt{--enable-lora},
\texttt{--max-lora-rank 64}); both servers run on the same GPU with the memory
fractions in Table~\ref{tab:serving}. Prefix caching is enabled on both
servers, and the cache is reset through the development endpoint
\texttt{/reset\_prefix\_cache} before each arm so that neither arm reuses the
other's KV cache. Requests are OpenAI-compatible chat completions with
temperature 0 and seed 0; $\tau^2$-bench and BFCL use the \texttt{xgrammar} guided
decoding backend and the \texttt{qwen3} reasoning parser.

\begin{table}[h]
\centering
\small
\caption{vLLM server settings per benchmark (target / draft on the same H100). A dash marks a setting
not used for that benchmark. On TaskBench at $k{=}5$ (Table~\ref{tab:model}) a round issues up to six
verification requests, so those beyond the fourth wait for a free slot.}
\label{tab:serving}
\begin{tabular}{lcccc}
\toprule
 & OpenAGI & TaskBench & $\tau^2$-bench & BFCL \\
\midrule
max model length (tokens) & 2{,}048 & 4{,}096 & 16{,}384 & 24{,}576 \\
GPU memory fraction, target / draft & 0.82 / 0.10 & 0.85 / 0.10 & 0.88 / 0.08 & 0.88 / 0.06 \\
max concurrent sequences, target / draft & 8 / 2 & 4 / 2 & 4 / 4 & 4 / 1 \\
guided decoding & -- & -- & xgrammar & xgrammar \\
reasoning parser & -- & -- & qwen3 & qwen3 \\
verification & \multicolumn{4}{c}{prefix-ordered, cancel} \\
\bottomrule
\end{tabular}
\end{table}

\paragraph{Common configurations of the size study.}
The rows of Table~\ref{tab:four} serve each benchmark's three drafts in one configuration that fits the 4B
draft beside the target: target / draft memory fractions 0.82 / 0.14 on OpenAGI and TaskBench, 0.83 / 0.13
with a 12{,}288-token context on $\tau^2$-bench, and 0.84 / 0.135 with a 15{,}360-token context on BFCL, all
with prefix-ordered cancellation; the Gemma target uses the same configurations.

\paragraph{The Gemma target.} Gemma-4-31B-it runs with FP8 weights on one H100: in bf16 it leaves a KV cache of
7{,}500 tokens at 85\% of the GPU, too few for $k{+}1$ concurrent verifications beside a draft. Its request
settings are those of the Qwen3 runs. Gemma departs from the prescribed action syntax more often than Qwen3
(Markdown fences, its native tool-call markup, keyword-argument calls on TaskBench, several calls in one
decision on $\tau^2$-bench); a surface-syntax normalization rewrites these forms into the prescribed one
before parsing, is applied identically to collection and to both arms, and every rewrite is counted (33 of
the 949 parsed actions of a TaskBench run).

\paragraph{Verification requests.}
The $k{+}1$ verification contexts of a round are issued as concurrent HTTP
requests to the target server. In the \emph{gather} mode of our earlier runs (no row of Table~\ref{tab:model} uses it), the round waits for all $k{+}1$ responses before
committing; the time spent waiting for responses that the first mismatch has
already made irrelevant is the idle term of App.~\ref{app:predictor}. In the \emph{prefix-ordered} mode used for every run we report, responses are consumed
in prefix order; as soon as context $i$ mismatches, the client closes the connections of
contexts $i{+}1,\dots,k$, which makes vLLM abort those requests, and the round
commits without waiting for them. The mode changes only when the round returns,
not what it commits.

\paragraph{Timing protocol.}
The test tasks of a shard are run first sequentially (target only) and then speculatively
on the same node and the same servers, with a cache reset in between; the
reported speedup is the pooled ratio of the summed whole-task wall-clocks.
OpenAGI was additionally run with prefix caching disabled on both servers as a
control (1.35\xx{} vs.\ 1.46\xx{} with caching, in the earlier gather-mode run). In Table~\ref{tab:four}, the
Qwen3 OpenAGI and TaskBench rows pool two independent passes over the test tasks (both passes of a task are
attached to it in the bootstrap); $\tau^2$-bench is run in two shards of 20 tasks and BFCL in three shards of
20, each sequential then speculative; the Gemma rows are single passes. The main runs of Table~\ref{tab:model}
were repeated once on $\tau^2$-bench and BFCL (App.~\ref{app:success}).


\section{Segment truncation by task length}
\label{app:length}
\begin{table}[h]
\centering
\caption{\textbf{Segment truncation by task length.} Tasks of the main runs bucketed by their sequential
length. $e$ (Eq.~\ref{eq:c}), the hit-rate sum $\sum_j q_j$ and $c$ are measured per bucket; the
prediction replays the bucket's sequential traces with the run-wide hit rates truncated at the task end,
so it carries the truncation effect only. Buckets hold 7--27 tasks, so individual ratios are noisy.}
\label{tab:length}
\small
\begin{tabular}{llrrrrrr}
\toprule
 & & & & & & \multicolumn{2}{c}{speedup} \\
\cmidrule(lr){7-8}
benchmark & steps & tasks & $e$ & $\sum_j q_j$ & $c$ & meas. & pred. \\
\midrule
OpenAGI   & 3        & 15 & 0.52 & 1.66 & 2.14 & 1.54\xx{} & 1.35\xx{} \\
OpenAGI   & 4        & 21 & 0.31 & 1.46 & 2.15 & 1.52\xx{} & 1.39\xx{} \\
OpenAGI   & 5        & 15 & 0.40 & 1.54 & 2.14 & 1.60\xx{} & 1.57\xx{} \\
OpenAGI   & $\ge 6$  &  9 & 0.30 & 1.10 & 1.80 & 1.39\xx{} & 1.55\xx{} \\
\midrule
TaskBench & $\le 4$  & 27 & 0.38 & 1.29 & 1.91 & 1.48\xx{} & 1.37\xx{} \\
TaskBench & 5        & 16 & 0.40 & 1.40 & 2.00 & 1.49\xx{} & 1.52\xx{} \\
TaskBench & 6        &  7 & 0.32 & 1.53 & 2.21 & 1.50\xx{} & 1.50\xx{} \\
TaskBench & $\ge 7$  & 10 & 0.12 & 0.77 & 1.65 & 1.14\xx{} & 1.52\xx{} \\
\midrule
$\tau^2$-bench   & $\le 9$  &  9 & 0.20 & 0.95 & 1.75 & 1.47\xx{} & 1.28\xx{} \\
$\tau^2$-bench   & 10--13   &  9 & 0.12 & 1.08 & 1.96 & 1.65\xx{} & 1.51\xx{} \\
$\tau^2$-bench   & 14--17   & 10 & 0.11 & 0.98 & 1.87 & 1.66\xx{} & 1.70\xx{} \\
$\tau^2$-bench   & $\ge 18$ & 12 & 0.07 & 0.96 & 1.89 & 1.55\xx{} & 1.72\xx{} \\
\bottomrule
\end{tabular}
\end{table}

Table~\ref{tab:length} buckets the tasks of the main runs by length. The truncation loss $e$ behaves as
Eq.~\ref{eq:c} says: it falls from 0.52 on three-step OpenAGI tasks to 0.30 on tasks of six or more
steps, and from 0.20 to 0.07 across $\tau^2$-bench's length range. The per-bucket speedup does not fall with
it, because task length and acceptance are not independent: short OpenAGI tasks are also the ones the
draft predicts best ($\sum_j q_j$ 1.66 against 1.10 on the longest tasks), and the longest TaskBench
tasks are where acceptance collapses ($\sum_j q_j$ 0.77, $c$ 1.65, 1.14\xx{}). A prediction that replays
each bucket with the run-wide hit rates captures the truncation effect but not this correlation, and
misses by up to 0.4 in the extreme buckets, whereas the run-level prediction is within 0.10. The
truncation term is therefore visible in the data, and the benchmarks' short tasks cost less than $e$
alone suggests because the draft is most often right on them.

\section{Transfer across benchmarks}
\label{app:transfer}
Table~\ref{tab:transfer} scores each trained 0.6B draft, and the untrained base, on the held-out contexts of
every benchmark with that benchmark's evaluator (Table~\ref{tab:offline}). All twenty cells come from one
job. The diagonal and the base row were measured again there; they differ from Table~\ref{tab:offline} by at
most 0.5 points, as repeated evaluations of the same draft do (the size study measured 75.4 and 74.7,
Table~\ref{tab:size}).

\begin{table}[h]
\centering
\caption{\textbf{Transfer across benchmarks.} Offline agreement (\%) of each 0.6B draft (rows) on the
held-out contexts of each benchmark (columns), scored as in Table~\ref{tab:offline} (BFCL: decision agreement
on the 99 development contexts).}
\label{tab:transfer}
\small
\begin{tabular}{lrrrr}
\toprule
 & \multicolumn{4}{c}{evaluated on} \\
\cmidrule(lr){2-5}
draft trained on & OpenAGI & TaskBench & $\tau^2$-bench & BFCL \\
\midrule
OpenAGI        & \textbf{75.5} & 29.6 & 8.8 & 3.0 \\
TaskBench      & 23.3 & \textbf{74.7} & 9.3 & 6.1 \\
$\tau^2$-bench & 10.8 & 21.7 & \textbf{63.7} & 1.0 \\
BFCL           & 12.0 & 23.1 & 8.0 & \textbf{66.7} \\
\midrule
untrained base & 13.3 & 25.3 & 8.0 & 2.0 \\
\bottomrule
\end{tabular}
\end{table}

\section{Online training details}
\label{app:online}
\paragraph{Stream and servers.} The training tasks of the benchmark are shuffled once with a fixed seed and
run in batches through the speculative system (target and draft on one GPU, the serving configuration of
App.~\ref{app:serving}). On OpenAGI each batch runs both arms. On $\tau^2$-bench the stream runs the
speculative arm only; the sequential arm on the same training tasks was run separately in the same serving
configuration. Fig.~\ref{fig:online}b reports, for both benchmarks, sequential time per step over speculative time
per committed step on each batch. Training runs on a second GPU while the servers wait between batches. The new adapter is
loaded into the running draft server by name through its runtime adapter endpoint, and the next batch proposes
with it. Every measurement on the stream is prequential: batch $b$ is never seen by the adapter that serves it.

\paragraph{Pairs.} On OpenAGI the commits of a batch are the (context, target action) pairs directly, with
duplicate (task, step) entries removed. $\tau^2$-bench records keep only hashes of the prompts, so the pairs are
rebuilt by replaying the committed actions in the benchmark environment, exactly as the collection pass
builds a step; every rebuilt prompt is checked against the hash the run recorded for the draft's request at
that step (898 pairs, no mismatch).

\paragraph{Recipe.} The adapter of every version is trained from the base with the recipe of
Table~\ref{tab:recipe}. Checkpoint selection uses a task-level 10\% carve of the buffer on OpenAGI and the
frozen 14 validation tasks on $\tau^2$-bench. The $\tau^2$-bench validation pairs were collected with the target sampling at
temperature 0.6, while the commits of the running system come from the verification setting (temperature 0);
the agreement of every $\tau^2$-bench draft in this paper, offline or online, is measured on the same validation
labels, so the curves are comparable, and the end-to-end acceptance is the measure under the verification
setting.

\paragraph{Continual variant.} Each version continues the previous adapter for two epochs on the new batch's
pairs plus an equal-size uniform sample of the earlier pairs, with the learning rate and rank of the recipe;
the first version is a fresh adapter on the first batch. The replay ratio and the number of epochs were not
tuned.

\begin{table}[tbp]
\centering\small\setlength{\tabcolsep}{8pt}
\caption{\textbf{Final online adapters on the test split, Qwen3 target.} Whole-task inference speedup
and task-bootstrap 95\% intervals, using trained 0.6B drafts with matched offline and online serving settings. Training time is reported separately.}
\label{tab:online_final_ci}
\begin{tabular}{lccc}
\toprule
Benchmark & Offline & Online: retrain & Online: continual \\
\midrule
OpenAGI & \shortstack{1.51\xx{}\\{\fontsize{8}{9}\selectfont [1.38, 1.64]}} & \shortstack{1.47\xx{}\\{\fontsize{8}{9}\selectfont [1.31, 1.63]}} & \shortstack{1.51\xx{}\\{\fontsize{8}{9}\selectfont [1.40, 1.64]}} \\
$\tau^2$-bench & \shortstack{1.61\xx{}\\{\fontsize{8}{9}\selectfont [1.44, 1.79]}} & \shortstack{1.61\xx{}\\{\fontsize{8}{9}\selectfont [1.43, 1.81]}} & \shortstack{1.63\xx{}\\{\fontsize{8}{9}\selectfont [1.46, 1.83]}} \\
\bottomrule
\end{tabular}
\end{table}


\paragraph{Endpoint.} The final adapter is measured on the test split with the protocol of
Table~\ref{tab:four} ($\tau^2$-bench: two shards of 20 tasks, each sequential then speculative, on separate server
pairs). Table~\ref{tab:online} lists the agreement of every version and the update times.

\begin{table}[h]
\centering
\caption{\textbf{Online training: agreement of each adapter version} (offline exact match on the held-out
contexts, \%). OpenAGI: 263 test contexts, one version per 25 training tasks; $\tau^2$-bench: 226 validation
contexts, one version per 10 tasks, so its 60 training tasks give six versions. Offline drafts: 75.0 (OpenAGI) and 63.7 ($\tau^2$-bench). Update time is the
time from the end of a batch to the loaded adapter (buffer build and training).}
\label{tab:online}
\small
\resizebox{\linewidth}{!}{\begin{tabular}{lrrrrrrrrrrr}
\toprule
version & 1 & 2 & 3 & 4 & 5 & 6 & 7 & 8 & 9 & 10 & update time \\
\midrule
OpenAGI, retrain on buffer & 53.0 & 58.8 & 67.0 & 69.6 & 73.3 & 74.7 & 73.4 & 72.8 & 74.1 & 74.0 & 27--73 s \\
OpenAGI, continual         & 48.7 & 61.6 & 65.7 & 69.8 & 63.5 & 69.2 & 70.8 & 72.5 & 75.9 & 74.9 & 22 s \\
$\tau^2$-bench, retrain on buffer & 33.6 & 57.5 & 58.4 & 58.4 & 64.6 & 66.4 & -- & -- & -- & -- & 3--18 min \\
$\tau^2$-bench, continual         & 28.3 & 57.1 & 54.9 & 59.7 & 63.3 & 63.3 & -- & -- & -- & -- & 3--6 min \\
\bottomrule
\end{tabular}}
\end{table}


\section{Task success and trajectory identity}
\label{app:success}
Table~\ref{tab:success} lists task outcomes for every configuration in Table~\ref{tab:four}. $\tau^2$-bench and BFCL use the benchmark's official success; OpenAGI and TaskBench have no
task-level success metric beyond termination, and the count is the number of runs that terminated. The
Qwen3 OpenAGI and TaskBench counts are over two runs of 63 and 60 tasks. The two passes of the Qwen3
$\tau^2$-bench main run gave 22 and 22 sequential against 21 and 20 speculative successes, and the two passes of the
BFCL main run 34 and 34 against 37 and 35. Across the 12 tool-benchmark size configurations and these four repeat comparisons,
official task success increases in 8 cases, decreases in 5, and is unchanged in 3.
Differences are at most three tasks per comparison. The 12 planning-benchmark comparisons instead
count normal termination: 3 decrease and 9 are unchanged.

Table~\ref{tab:parity} compares whole trajectories. The target run alone is reproducible: every pair of
target-only runs of the same configuration commits identical decisions on every task. Speculative runs are not,
mainly because verification decodes several contexts together and the engine's floating-point results depend on which
requests share a batch. On the planning benchmarks, where a target step takes 1--2\,s, 83--96\% of the
speculative trajectories are identical to the target-only one. On the tool benchmarks, where a step reasons for
10--12\,s, a small numerical difference changes some decision on most tasks, and the trajectory differs from
then on; two speculative runs also differ from each other. Every one of the 938 speculative task runs in
Table~\ref{tab:parity} passed the harness's commit audit, which replays the task and checks that each committed
action is the target's output on the context in which it was committed.

\paragraph{Serialized verification.} To separate batching from other effects, we reran speculation with the
target server limited to one sequence at a time, so that every verification request is decoded alone, as in the
target-only run (TaskBench: 60 tasks, both arms; $\tau^2$-bench: the first 20 test tasks, compared with the
target-only arm of the main run). The limit does not change the target-only run (60 of 60 TaskBench tasks
identical). The speculative trajectories then match the target-only ones on 57 of 60 TaskBench tasks, against 50
of 60 with concurrent verification, and on 8 of 20 $\tau^2$-bench tasks, against 0 of 20; every commit passes the
audit. Batching thus accounts for most of the divergence but not for all of it; we did not identify the other
source within the serving engine.

\begin{table}[!ht]
\centering
\caption{\textbf{Task outcomes per arm} for the configurations of Table~\ref{tab:four}.
Counts denote normal termination on OpenAGI and TaskBench, and official task success on
$\tau^2$-bench and BFCL. The sequential count is the same across the three draft-size comparisons.}
\label{tab:success}
\small
\begin{tabular}{lrrrrrrrr}
\toprule
 & \multicolumn{4}{c}{Qwen3-32B target} & \multicolumn{4}{c}{Gemma-4-31B target} \\
\cmidrule(lr){2-5}\cmidrule(lr){6-9}
 & & \multicolumn{3}{c}{speculative} & & \multicolumn{3}{c}{speculative} \\
\cmidrule(lr){3-5}\cmidrule(lr){7-9}
benchmark & sequential & 0.6B & 1.7B & 4B & sequential & 0.6B & 1.7B & 4B \\
\midrule
OpenAGI         & 126 & 126 & 126 & 126 & 63 & 63 & 63 & 63 \\
TaskBench       & 116 & 114 & 115 & 114 & 60 & 60 & 60 & 60 \\
$\tau^2$-bench  &  22 &  23 &  24 &  21 & 39 & 40 & 40 & 40 \\
BFCL multi-turn &  34 &  34 &  36 &  34 & 48 & 47 & 47 & 48 \\
\bottomrule
\end{tabular}

\vspace{1.5em}
\caption{\textbf{Trajectory identity.} Share of tasks (\%) on which two runs of the same configuration commit
identical decisions, with the free text of replies and summaries ignored; comparing that text as well gives
3.3\% on BFCL and 0\% on $\tau^2$-bench for target-only against speculative. OpenAGI and TaskBench: the six
runs of the draft-size study (three drafts, two repeats); $\tau^2$-bench and BFCL: the two passes of the main
configuration.}
\label{tab:parity}
\small
\begin{tabular}{lrrr}
\toprule
benchmark & two target-only runs & target-only vs.\ speculative & two speculative runs \\
\midrule
OpenAGI        & 100 & 95.5 & 96.5 \\
TaskBench      & 100 & 82.8 & 81.7 \\
$\tau^2$-bench & 100 &  1.3 & 45.0 \\
BFCL           & 100 & 34.2 & 60.0 \\
\bottomrule
\end{tabular}
\end{table}
